\section{Problem Statement}
\label{sec:problem_statement}
Ethereum Layer-1 is constrained by limited throughput (historically processing 15-25 TPS) and high fee volatility \cite{buterin2014ethereum}. Zero-knowledge rollups (ZK rollups) have emerged as a promising Layer-2 (L2) solution to address these scalability issues. By processing thousands of transactions off-chain and submitting compact validity proofs to L1, ZK rollups can theoretically achieve significantly higher throughput while inheriting Ethereum’s robust security model \cite{tran2023implementation}. 

However, real-world deployments of ZK rollups reveal critical bottlenecks that limit their practical effectiveness. Practical ZK-Rollup performance is dictated by several interacting and often competing factors. These include batch formation limits, sequencer scheduling algorithms, data availability footprints, Layer-1 submission costs, and compute-intensive proof-generation overhead \cite{habib2025bottlenecks,chaliasos2024analyzing}. Importantly, these factors do not affect performance independently. For instance, increasing the maximum batch size reduces the amortized Layer-1 data cost per transaction but simultaneously increases the queueing delay (latency) for early transactions waiting in the pool for the batch to seal.

Currently, observing these trade-offs requires either complex theoretical modeling or passive analysis of production mainnet systems (like zkSync Era or Starknet). Production systems are highly optimized for reliability, full EVM compatibility, and end-user security rather than for scientific experimentation. They are difficult to modify, and their real-world workloads cannot be perfectly replicated to test different configurations under identical conditions.

Therefore, the core research problem addressed in this report is the lack of a compact, modular, and reproducible experimental framework designed specifically to isolate how batching, sequencing, proof-delay, and data availability choices affect ZK-Rollup throughput, latency, finalization time, and cost behavior under strictly comparable, controlled synthetic workloads.
